% preamble-exam-zh.tex — 极简讲解页共享 preamble
% 目标：复刻「题目独立、提示轻框、路线箭头、主体双栏、答案轻收束」的讲义风格。

\documentclass{exam-zh}

% ---------- 基础配置 ----------
\examsetup{
  page/size = a4paper,
  page/show-head = false,
  page/show-foot = false,
  sealline/show = false,
  font = none,
  math-font = none,
}

% ---------- 包 ----------
\usepackage{geometry}
\geometry{top=10mm,bottom=14mm,left=15mm,right=13mm}
\AtBeginDocument{\newgeometry{top=10mm,bottom=14mm,left=15mm,right=13mm}}

\usepackage{xcolor}
\usepackage{needspace}
\usepackage{enumitem}
\usepackage{array}
\usepackage{tabularx}
\usepackage{tcolorbox}
\tcbuselibrary{skins,breakable}

% ---------- 打印/屏幕颜色切换 ----------
% 用法：\documentclass[monochrome]{exam-zh} 可降级为灰度。
% 注意：不要使用 \if@xxx 放在 makeatother 外部；这里改成普通条件名。
\newif\ifeduprintcolor
\eduprintcolortrue
\makeatletter
\@ifclasswith{exam-zh}{monochrome}{\eduprintcolorfalse}{}
\makeatother

% ---------- 语义颜色 ----------
\ifeduprintcolor
  \definecolor{edu-blue}{RGB}{31,62,126}
  \definecolor{edu-blue-soft}{RGB}{232,238,250}
  \definecolor{edu-blue-bg}{RGB}{248,250,254}
  \definecolor{edu-ink}{RGB}{30,35,45}
  \definecolor{edu-muted}{RGB}{118,124,136}
  \definecolor{edu-line}{RGB}{209,218,235}
  \definecolor{edu-soft-bg}{RGB}{250,251,253}
  \definecolor{edu-red}{RGB}{168,58,58}
  \definecolor{edu-red-bg}{RGB}{255,247,245}
\else
  \definecolor{edu-blue}{gray}{0.15}
  \definecolor{edu-blue-soft}{gray}{0.90}
  \definecolor{edu-blue-bg}{gray}{0.98}
  \definecolor{edu-ink}{gray}{0.10}
  \definecolor{edu-muted}{gray}{0.45}
  \definecolor{edu-line}{gray}{0.82}
  \definecolor{edu-soft-bg}{gray}{0.97}
  \definecolor{edu-red}{gray}{0.28}
  \definecolor{edu-red-bg}{gray}{0.97}
\fi

% ---------- 全局排版 ----------
\setlength{\parindent}{0pt}
\setlength{\parskip}{0pt}
\linespread{1.16}
\renewcommand{\arraystretch}{1.15}

% ctex 通常提供 \kaishu；这里用它做教师旁白感。
\newcommand{\edukai}{\kaishu}

% ---------- 顶部元信息 ----------
\newcommand{\eduTopBar}[2]{%
  \noindent
  \begin{tabularx}{\linewidth}{@{}Xr@{}}
    {\color{edu-blue}\bfseries #1} &
    {\color{edu-blue}\bfseries #2}
  \end{tabularx}
  \vspace{8mm}
}

% ---------- 轻标题体系 ----------
% 只有真正的大结构才用它；不要给提示/路线滥用标题。
\newcommand{\eduSectionTitle}[1]{%
  \vspace{2mm}
  {\color{edu-blue}\bfseries\large #1}\par
  \vspace{3mm}
}

\newcommand{\eduSubTitle}[1]{%
  \vspace{1mm}
  {\color{edu-blue}\bfseries #1}\par
  \vspace{1mm}
}

% ---------- 小问题干：复用 exam (1)(2)(3) 格式 ----------
\newenvironment{eduSubQuestionItem}[1]{%
  \vspace{2mm}
  \par\noindent
  \hangindent=8mm\hangafter=1
  {\color{edu-blue}\bfseries #1}\enspace
  \begingroup
  \color{edu-ink}
}{%
  \endgroup
  \par
  \vspace{3mm}
}

% ---------- 辅助信息条（解题流程/关键想法/思考）楷体 ----------
\newenvironment{eduAuxItem}[1]{%
  \vspace{2mm}
  \par\noindent
  \hangindent=8mm\hangafter=1
  {\color{edu-blue}\edukai $\triangleright$~#1}\enspace
  \begingroup
  \color{edu-ink}
  \edukai
}{%
  \endgroup
  \par
  \vspace{4mm}
}

\newcommand{\eduColumnTitle}[2]{%
  {\color{edu-blue}\bfseries #1\hspace{1.5mm}#2}\par
  \vspace{2.5mm}
}

% ---------- 图标：不用外部 icon 字体，保证可移植 ----------
\newcommand{\eduIconBulb}{\(\circ\)}
\newcommand{\eduIconBook}{\(\square\)}
\newcommand{\eduIconCheck}{\(\checkmark\)}
\newcommand{\eduIconPin}{\(\diamond\)}
\newcommand{\eduIconNote}{\(\triangleright\)}

% ---------- 题目区 ----------
% 题目区是页面入口：弱标签 + 一条长蓝线 + 大题干。
\newenvironment{eduProblemBlock}[1][题目]{%
  \needspace{8\baselineskip}
  {\small\color{edu-muted}#1}\par
  \vspace{1.5mm}
  {\color{edu-blue}\hrule height 0.55pt}
  \vspace{4.5mm}
  \begingroup
  \color{edu-ink}
  \large
  \setlength{\parskip}{2.2mm}
  \linespread{1.20}\selectfont
}{%
  \par
  \endgroup
  \vspace{6mm}
}

\newenvironment{eduSubQuestions}{%
  \begin{enumerate}[
    label=(\arabic*),
    leftmargin=8mm,
    itemsep=2.0mm,
    topsep=2.0mm,
    parsep=0pt
  ]
}{%
  \end{enumerate}
}

% ---------- 读题提示：轻框，不做同级标题 ----------
\newtcolorbox{eduReadTipBox}{
  enhanced,
  breakable,
  colback=white,
  colframe=edu-blue!45,
  boxrule=0.45pt,
  arc=1.6mm,
  left=10mm,
  right=4mm,
  top=5mm,
  bottom=5mm,
  before skip=1mm,
  after skip=7mm,
  fontupper=\edukai\normalsize,
  colupper=edu-ink,
  overlay={
    \node[anchor=north west, text=edu-blue] at ([xshift=3.2mm,yshift=-3.2mm]frame.north west)
    {\small\eduIconNote};
  }
}

\newenvironment{eduTipList}{%
  \begin{itemize}[
    label=\textbullet,
    leftmargin=5mm,
    itemsep=1.2mm,
    topsep=0pt,
    parsep=0pt
  ]
}{%
  \end{itemize}
}

% ---------- 解题路线：虚线框 + 文字箭头 ----------
\newenvironment{eduRouteFlow}{%
  \vspace{1mm}
  \begin{center}
  \normalsize
}{%
  \end{center}
  \vspace{7mm}
}
\newcommand{\eduRouteStep}[1]{%
  {\color{edu-ink}#1}%
}
\newcommand{\eduRouteArrow}{%
  \;$\boldsymbol{\to}$\;%
}

% ---------- 主体讲解：只在这里使用双栏 ----------
% 不用 tabularx 包整块，避免模板生成后出现 tabularx body 扫描问题。
\newenvironment{eduDualExplain}{%
  \needspace{10\baselineskip}
  \par\noindent
}{%
  \par\vspace{5mm}
}

\newenvironment{eduExplainLeft}{%
  \begin{minipage}[t]{0.30\linewidth}
  \normalsize
  \color{edu-ink}
}{%
  \end{minipage}
}

\newcommand{\eduColumnDivider}{%
  \hspace{4mm}{\color{edu-blue}\vrule width 0.45pt}\hspace{7mm}%
}

\newenvironment{eduExplainRight}{%
  \begin{minipage}[t]{0.58\linewidth}
  \normalsize
  \color{edu-ink}
}{%
  \end{minipage}
}

\newenvironment{eduNumberList}{%
  \begin{enumerate}[
    label=\arabic*.,
    leftmargin=6mm,
    itemsep=4.8mm,
    topsep=0pt,
    parsep=0pt
  ]
}{%
  \end{enumerate}
}

\newenvironment{eduBulletList}{%
  \begin{itemize}[
    label=\textbullet,
    leftmargin=5mm,
    itemsep=1.2mm,
    topsep=0pt,
    parsep=0pt
  ]
}{%
  \end{itemize}
}

% ---------- 底部双栏：答案 + 方法提醒 ----------
\newenvironment{eduBottomGrid}{%
  \vspace{1mm}
  \par\noindent
}{%
  \par\vspace{1mm}
}

\newenvironment{eduSummaryLeft}{%
  \begin{minipage}[t]{0.30\linewidth}
}{%
  \end{minipage}
}

\newcommand{\eduSummaryDivider}{%
  \hspace{4mm}{\color{edu-line}\vrule width 0.35pt}\hspace{7mm}%
}

\newenvironment{eduSummaryRight}{%
  \begin{minipage}[t]{0.58\linewidth}
}{%
  \end{minipage}
}

% ---------- 答案/方法提醒：轻量收束 ----------
\newtcolorbox{eduSummaryBlock}[2]{
  enhanced,
  breakable,
  colback=white,
  colframe=white,
  boxrule=0pt,
  sharp corners,
  left=0mm,
  right=0mm,
  top=1mm,
  bottom=1mm,
  before skip=1mm,
  after skip=1mm,
  title={\color{edu-blue}\bfseries #1\hspace{1.5mm}#2},
  coltitle=edu-blue,
  attach title to upper={\par\vspace{2mm}},
  fontupper=\normalsize
}

\newenvironment{eduPlainList}{%
  \begin{enumerate}[
    label=(\arabic*),
    leftmargin=7mm,
    itemsep=1.2mm,
    topsep=0pt,
    parsep=0pt
  ]
}{%
  \end{enumerate}
}

% ---------- 兼容旧 block 的轻量盒子 ----------
\newtcolorbox{eduSoftBox}{
  enhanced,
  breakable,
  colback=edu-soft-bg,
  colframe=edu-line,
  boxrule=0.35pt,
  arc=1mm,
  left=3mm,
  right=3mm,
  top=2mm,
  bottom=2mm,
  before skip=1mm,
  after skip=2mm,
  fontupper=\small
}

\newtcolorbox{eduMistakeBox}{
  enhanced,
  breakable,
  colback=edu-red-bg,
  colframe=edu-red!40,
  boxrule=0.35pt,
  arc=1mm,
  left=3mm,
  right=3mm,
  top=2.5mm,
  bottom=2.5mm,
  before skip=1mm,
  after skip=2mm,
  fontupper=\normalsize
}

\newtcolorbox{eduLegacySolution}{
  enhanced,
  breakable,
  colback=edu-soft-bg,
  colframe=edu-line,
  boxrule=0.35pt,
  arc=1mm,
  left=3mm,
  right=3mm,
  top=2mm,
  bottom=2mm,
  before skip=-2mm,
  after skip=4mm,
  fontupper=\small
}

\newtcolorbox{eduTeacherNote}{
  enhanced,
  breakable,
  colback=edu-soft-bg,
  colframe=edu-line,
  boxrule=0pt,
  borderline west={1.5pt}{0pt}{edu-muted},
  left=2.5mm,
  right=2.5mm,
  top=2mm,
  bottom=2mm,
  before skip=1mm,
  after skip=2mm,
  fontupper=\footnotesize,
  colupper=edu-muted
}

% ---------- 答题区 ----------
\newcommand{\answerarea}[1][40mm]{%
  \vspace{3mm}%
  \begin{tcolorbox}[
    enhanced,
    breakable,
    colback=white,
    colframe=edu-line,
    boxrule=0.55pt,
    arc=0.8mm,
    height=#1,
    valign=top,
  ]
  \end{tcolorbox}%
}

\examsetup{
  question/show-answer = true,
  fillin/show-answer = true,
  solution/show-solution = show-stay,
}

\begin{document}


\begin{eduProblemBlock}[题目]
已知一次函数 $y=kx+b$ 的图像经过点 $A(1,3)$ 和点 $B(-2,-6)$。
\begin{enumerate}[label=(\arabic*)]
  \item 求这个一次函数的解析式；
  \item 判断点 $C(3,9)$ 是否在这个函数的图像上；
  \item 求该函数图像与 $x$ 轴、$y$ 轴的交点坐标，并计算图像与坐标轴围成的三角形面积。
\end{enumerate}

\end{eduProblemBlock}









\begin{eduAuxItem}{解题流程}
两点代入 $y=kx+b$，列方程组$\;\to\;$解方程组求 $k$、$b$$\;\to\;$代入检验点 $C$$\;\to\;$求坐标轴交点，算面积\end{eduAuxItem}




\begin{eduAuxItem}{关键想法}
一次函数有两个未知系数 $k$ 和 $b$。"图像经过一个点"意味着这个点的坐标满足解析式。
两个点 = 两个方程，刚好解出两个未知数。
这是待定系数法最基本的应用。
\end{eduAuxItem}




\needspace{12\baselineskip}
\begin{eduSubQuestionItem}{(1)}
求这个一次函数的解析式；\end{eduSubQuestionItem}

\begin{eduDualExplain}
  \begin{eduExplainLeft}
    \eduColumnTitle{\eduIconBulb}{思考引导}
    \begin{eduNumberList}
      \item 把 $A(1,3)$ 代入 $y=kx+b$，你能得到什么方程？      \item 两个方程，用什么方法消元最快？试试两式直接相减。    \end{eduNumberList}
  \end{eduExplainLeft}
  \eduColumnDivider
  \begin{eduExplainRight}
    \eduColumnTitle{\eduIconBook}{规范讲解}
    \begin{eduNumberList}
      \item 代入 $A(1,3)$：$3 = k + b$ \quad \textcircled{1}      \item 代入 $B(-2,-6)$：$-6 = -2k + b$ \quad \textcircled{2}      \item \textcircled{1}$-$\textcircled{2}：$9 = 3k$，所以 $k = 3$      \item 代回 \textcircled{1}：$3 = 3 + b$，所以 $b = 0$      \item 解析式：$y = 3x$    \end{eduNumberList}

    \vspace{1mm}
    \eduSubTitle{注意}
    \begin{eduBulletList}
      \item $b = 0$！这条直线过原点——后面的第（3）问会用到这个事实。    \end{eduBulletList}
  \end{eduExplainRight}
\end{eduDualExplain}




\needspace{12\baselineskip}
\begin{eduSubQuestionItem}{(2)}
判断点 $C(3,9)$ 是否在这个函数的图像上；\end{eduSubQuestionItem}

\begin{eduDualExplain}
  \begin{eduExplainLeft}
    \eduColumnTitle{\eduIconBulb}{思考引导}
    \begin{eduNumberList}
      \item 判断一个点是否在函数图像上，怎么做？      \item 把 $x$ 坐标代入解析式，看算出来的 $y$ 是否一致。    \end{eduNumberList}
  \end{eduExplainLeft}
  \eduColumnDivider
  \begin{eduExplainRight}
    \eduColumnTitle{\eduIconBook}{规范讲解}
    \begin{eduNumberList}
      \item 将 $x = 3$ 代入 $y = 3x$：$y = 3 \times 3 = 9$      \item 点 $C$ 的纵坐标也是 $9$，所以 $C(3,9)$ 在图像上。    \end{eduNumberList}

    \vspace{1mm}
    \eduSubTitle{方法}
    \begin{eduBulletList}
      \item 代进去，算出来，比一比。判断「是否在图像上」只有这一招。    \end{eduBulletList}
  \end{eduExplainRight}
\end{eduDualExplain}




\needspace{12\baselineskip}
\begin{eduSubQuestionItem}{(3)}
求该函数图像与 $x$ 轴、$y$ 轴的交点坐标，并计算图像与坐标轴围成的三角形面积。\end{eduSubQuestionItem}

\begin{eduDualExplain}
  \begin{eduExplainLeft}
    \eduColumnTitle{\eduIconBulb}{思考引导}
    \begin{eduNumberList}
      \item 与 $x$ 轴交点：令 $y = 0$，解出 $x$。      \item 与 $y$ 轴交点：令 $x = 0$，解出 $y$。      \item \textbf{等等！}$b = 0$ 时直线过原点，两个交点会怎样？    \end{eduNumberList}
  \end{eduExplainLeft}
  \eduColumnDivider
  \begin{eduExplainRight}
    \eduColumnTitle{\eduIconBook}{规范讲解}
    \begin{eduNumberList}
      \item 令 $y = 0$：$3x = 0$，$x = 0$，交点为 $(0,0)$      \item 令 $x = 0$：$y = 3 \times 0 = 0$，交点为 $(0,0)$      \item 两个交点\textbf{重合于原点}，无法围成三角形      \item 面积 $S = 0$    \end{eduNumberList}

    \vspace{1mm}
    \eduSubTitle{关键}
    \begin{eduBulletList}
      \item 正比例函数 $y = kx$ 过原点，与两轴只交于一个点，面积必然为 $0$。      \item 只有 $b \neq 0$ 时，直线与两轴才有两个不同交点，才能围成三角形。    \end{eduBulletList}
  \end{eduExplainRight}
\end{eduDualExplain}




\begin{eduBottomGrid}
  \begin{eduSummaryLeft}
    \begin{eduSummaryBlock}{\eduIconCheck}{答案}
    \begin{eduPlainList}
      \item (1) $y = 3x$      \item (2) $C(3,9)$ 在图像上      \item (3) 与两轴交点均为 $(0,0)$，面积为 $0$    \end{eduPlainList}
    \end{eduSummaryBlock}
  \end{eduSummaryLeft}
  \eduSummaryDivider
  \begin{eduSummaryRight}
    \begin{eduSummaryBlock}{\eduIconPin}{方法提醒}
    \begin{eduBulletList}
      \item 待定系数法四步：代入 → 列方程 → 消元 → 回代      \item 每次求出参数后先看有没有特殊情况（如 $b = 0$）    \end{eduBulletList}
    \end{eduSummaryBlock}
  \end{eduSummaryRight}
\end{eduBottomGrid}




\begin{eduMistakeBox}
\eduSubTitle{套面积公式不看条件}直线与坐标轴围成三角形面积公式 $S = \frac{1}{2}|x_0| \cdot |y_0|$ 的前提是：
两个交点不重合。本题 $b = 0$，两个交点都是原点，公式虽然代入也得 $0$，
但如果换一题（比如两交点确实重合但你没注意到），思维就错了。
**养成习惯：先判断两个交点是否不同，再套公式。**
\end{eduMistakeBox}




\begin{eduMistakeBox}
\eduSubTitle{代入负号出错}将 $B(-2, -6)$ 代入 $y = kx + b$ 时，容易写成 $-6 = k \cdot (-2) + b$ 后
漏掉负号变成 $-6 = 2k + b$。正确应为 $-6 = -2k + b$。
\end{eduMistakeBox}




\begin{eduMistakeBox}
\eduSubTitle{求出 b=0 却没有反应}$b = 0$ 不是普通的数值结果——它意味着这条直线过原点（正比例函数）。
每次求出 $b = 0$，都应该立刻标注"过原点"，这会影响后续所有与坐标轴有关的计算。
\end{eduMistakeBox}




\begin{eduAuxItem}{思考}
如果把点 $B$ 改为 $(-2, -4)$，求出来的 $b$ 还等于 $0$ 吗？围成的面积会是多少？\end{eduAuxItem}




\begin{eduAuxItem}{思考}
一次函数 $y = kx + b$ 在什么条件下与坐标轴围成的三角形面积最大？（提示：固定面积 $S = 6$，你能写出几条满足条件的直线？）\end{eduAuxItem}




\begin{eduAuxItem}{思考}
如果一个一次函数与两轴围成面积为 $12$ 的三角形，且过点 $(1, 3)$，这个函数的解析式是什么？（注意可能有两组解）\end{eduAuxItem}



\clearpage

\begin{eduAuxItem}{关键想法}
本题核心陷阱在第（3）问的退化情形。
学生画像：初二中等生，能做待定系数法但对 b=0 没有警觉。
教学节奏：前两问让学生自己完成，第（3）问作为课堂讨论重点。
\end{eduAuxItem}


\begin{eduTeacherNote}
  \textbf{训练目标：}待定系数法 + b=0 退化识别\par
  \textbf{预期卡点：}第(3)问机械套面积公式，不判断两交点是否重合
\end{eduTeacherNote}


\end{document}
